\section{Introduction}
\label{sec:introduction}

Urban building energy modeling (UBEM) and city-scale digital twins require thousands of building attributes---height, footprint area, use type, thermal zoning proxies, and consumption indicators---assembled from heterogeneous sources including OpenStreetMap (OSM), census statistics, LiDAR rasters, utility networks, CityGML databases, and simulation outputs.
At city scale, \emph{ground-truth validation of every attribute is infeasible}.
Practitioners therefore operate on blended measured, inferred, and default values, yet most pipelines expose a \emph{single number per feature} without lineage or uncertainty.

In our prior CIM Wizard system (Paper~1), feature calculators expose multiple methods per attribute with priority-ordered fallback: the first successful method wins and only one value is stored for the baseline scenario.
Provenance is implicit and alternative methods are never computed.
This design optimises execution speed but hides methodological disagreement and prevents confidence-aware downstream UBEM.

This paper presents \textbf{CIM Wizard 2.0}, an architecture that makes urban modeling outputs \emph{auditable} when full accuracy cannot be established.
We pursue six engineering objectives (O1--O6) and six testable hypotheses (H1--H6):
\begin{enumerate}[label=\textbf{O\arabic*:}, leftmargin=*]
  \item A \textbf{semi data warehouse} registering heterogeneous datasources with STAC-aligned metadata.
  \item A \textbf{three-layer UBEM ontology stack} virtualised over PostgreSQL via Ontop OBDA.
  \item A \textbf{CIM-scoped text-to-SQL benchmark} for PostGIS multi-schema queries.
  \item \textbf{Fine-tuned spatial-SQL LLMs} validated with execution-centric metrics.
  \item A \textbf{multi-agent pipeline} automating VKG/OBDA mapping generation.
  \item \textbf{Integration} with multi-method execution, provenance records, and confidence scoring.
\end{enumerate}

Our central thesis is that \emph{structured provenance, multi-method comparison, and semantic lineage querying} provide a practical substitute for exhaustive ground-truth validation---supporting sensitivity analysis and confidence scoring until better measurements become available.

\subsection{Contributions}

\begin{itemize}[leftmargin=*]
  \item \textbf{O1/H1:} CIM semi data warehouse design (\texttt{datalake/}) with STAC registry, ingest API, and datasource versioning for provenance back-links.
  \item \textbf{O2/H2:} Layered UBEM ontology (meta / upper / domain) with alignment across CityGML, BOT, SOSA, and GeoSPARQL; seed OBDA mappings over CityDB and energy extensions.
  \item \textbf{O3/H3:} Taxonomy-structured spatial-SQL benchmark (\texttt{ai4db/}: rule-based Stage~1 + LLM-augmented Stage~3).
  \item \textbf{O4/H4:} Fine-tuned Qwen/Llama models and unified evaluator (\texttt{txt2ssql/}, \texttt{assist\_cim/}) reporting EM, EX, SC, and EA.
  \item \textbf{O5/H5:} Multi-agent VKG construction combining LLM4VKG mapping patterns, fine-tuned source-SQL generation, and template-based RDF targets.
  \item \textbf{O6/H6:} Provenance-aware multi-method pipeline with ordinal confidence levels and UBEM stock-sensitivity evaluation.
\end{itemize}

\subsection{Paper organisation}

Section~\ref{sec:related} reviews prior work per objective.
Section~\ref{sec:architecture} presents the system architecture.
Section~\ref{sec:methodology} describes the five-phase methodology.
Section~\ref{sec:evaluation} formalises hypotheses and metrics.
Section~\ref{sec:casestudy} reports the case study.
Section~\ref{sec:conclusion} concludes.
